\section{A Data-Driven Approach to Calibrating POFUL}

%In this paper, we study Thompson Sampling (TS) for the stochastic linear bandit problem. TS is a popular policy to use in practice partly due to its strong empirical performance. In addition, it is known that TS achieves minimax Bayesian regret. However, its worst-case regret suffers from an extra multiplicative factor of order $\sqrt{d}$, which is recently shown to be unavoidable. Motivated by this discrepancy between the worst-case versus typical performance of TS, we derive an instance-dependent regret bound for TS that can be calculated efficiently in a ``data-driven" manner through a geometric approach. This result enables us to detect and ``course-correct" if TS performs poorly in any specific problem instance. We show a variant of TS based on our approach enjoys the minimax optimal worst-case regret of $\tilde{\cO}(d\sqrt{T})$ while keeping most of the properties of TS. Simulation results are presented to validate our insights and to compare our algorithm with the baselines.





In this section, we present the main contribution of this work which provides a data-driven approach to calibrating POFUL. We'll focus on LinTS and Greedy as they lack theoretical results of frequentist regret bound. In these settings, we have $\tau_t = 0$ and hence $\mu_t  \coloneqq  (1+\alpha_t)(1+\iota_t )$. Note that $\iota_t$ is the inflation parameter that can be controlled by a decision-maker, the essential point is to find a computable, non-trivial upper bound for the sensitivity ratio $\alpha_t$, which turns into an upper bound for the regret ratio $\mu_t$ that's deeply related to the frequentist regret of POFUL.

 In the following, we first construct such upper bounds $\{\hat{\alpha}_t\}_{t\in[T]}$. Our strategy capitalizes on geometric insights related to the properties of the confidence ellipsoids. Subsequently, we show how these bounds can be utilized to calibrate POFUL. Specifically, we propose the TS-MR and Greedy-MR algorithms that attain a provable minimax-optimal frequentist regret up to constant factors, and yet keep the majority of their original properties such as computational efficiency and low regret in most instances.














%todo: add a paragraph about the meaning of $V_t$, add a confidence ellipsoid figure for illustrating 




 




\subsection{Upper Bounding the Sensitivity Ratio}
\label{sec:bound_sst_ratio} 

We consider two scenarios: (1) Continuous action space and (2) Discrete action space. For sake of a better illustration, we consider 
 $\cX_t \equiv \cS_{d-1}$ for the first scenario, which is a standard example of continuous action space, and is the same as the setting considered in \cite{abeille2017linear}. The analysis of the second scenario is deferred to Appendix~\ref{sec:discrete_actions}.
 
 %For the second scenario,  we consider uniformly distributed actions $\cX_t \overset{i.i.d.}{\sim} \text{Unif}(\cS_{d-1})$.



\paragraph{Continuous action space.} 
% In the following, we present our main result that lower bounds sensitivity ratios for continuous action spaces $\cX_t = \cS_{d-1}$ for all $t \in [T]$.  
Note that in this case, the optimal action $x_t^\star(\theta )  \coloneqq \argmax_{x\in \cX_t} \dotp{x}{\theta}$ takes a closed form $ x^\star_t(\theta ) = \theta/\|\theta\|$, we only need to study the property of the confidence ellipsoids of $\theta$, namely, $\cE^{\RLS}_t$ and $\cE^{\PVT}_t$. A direct idea is to consider respectively the smallest and largest value of $\|x^\star_t(\theta)\|_{V_t^{-1}}$ for $\theta$ in the ellipsoids. Specifically, we have 
  \# \label{eq:alpha_hat_raw}
 \alpha_t \leq \frac{\sup _{\theta \in \cE_t^{\RLS}}\|x^\star_t(\theta)\|_{V_t^{-1}}}{ \inf _{\theta \in \cE_t^{\PVT}}\|x^\star_t(\theta)\|_{V_t^{-1}}}.
 \#
 At first glance, one might think that \eqref{eq:alpha_hat_raw} is too conservative to provide useful calibration for POFUL. However, it turns out that the lower bound is pretty tight in practice. As is illustrated in Figure~\ref{fig:ellipsoid1}. The intuition is that, when POFUL has implemented sufficient exploration so that the confidence ellipsoid   $\cE_t$ is small enough, the potentially optimal actions concentrate in a small cap $\cC_t$ on the unit hypersphere $\cS_{d-1}$. All actions within $\cC_t$ roughly point to the same direction and thus have similar $V_t^{-1}$-norm. The following main theorem of this paper derives a upper bound of $\alpha_t$ based on this idea. We sketch the proof in Appendix~\ref{sec:proof_sketch_alpha}.
 \begin{figure}[htp]
	\centering
%	\subfigure{
		\includegraphics[width=.25\textwidth]{figures/ellipsoid_4} 
	%}
	\caption{Illustration of the set $\cC_t$ of potentially optimal actions in $\RR^2$.}
	\label{fig:ellipsoid1}
\end{figure}

The following assumption simplifies the proof, and we remark that it holds almost surely. 
\begin{assumption}
\label{assump:distinct_eigens}
	For all $t$, $V_t$ has distinct eigenvalues, i.e. $\lambda_1( V_t)> \lambda_2( V_t) \dots > \lambda_d( V_t) \geq 0$. 
\end{assumption}

\begin{theorem}\label{thm:alpha}

	Under Assumpsion~\ref{assump:distinct_eigens},  define $\beta^{\PVT}_t = \iota_t \beta^{RLS}_t$ and
					\$
	\Phi_t &=   \begin{cases}
				( \lambda_{1}^{-1} + \lambda_{d}^{-1}   - m_t\lambda_{1}^{-1}\lambda_{d}^{-1} )^{\frac{1}{2}}, &\ \ \ \ \textit{if}\  \|\hat{\theta}_t\|_{V_t} \geq \beta^{RLS}_t	\\
				\lambda_d^{-\frac{1}{2}}(V_t), & \ \ \ \ \textit{if}\  \|\hat{\theta}_t\|_{V_t} <  \beta^{RLS}_t \end{cases},\\
	\Psi_t &=   \begin{cases}
				(\lambda_{k}^{-1} +  \lambda_{k+1}^{-1}   - M_t\lambda_{k}^{-1}\lambda_{k+1}^{-1} )^{\frac{1}{2}}, &\textit{if}\  \|\hat{\theta}_t\|_{V_t} \geq \beta^{\PVT}_t	\\
				\lambda_1^{-\frac{1}{2}}(V_t), & \textit{if}\  \|\hat{\theta}_t\|_{V_t} <  \beta^{\PVT}_t\end{cases},
				\$ 
				
				

				
				
				
 where we let
	\$
	 \quad m_t = (\|\hat{\theta}_t\|_{V_t}^2  - (\beta^{RLS}_t)^2)/(\|\hat{\theta}_t\| +  \beta^{RLS}_t/\lambda_d(V_t) )^2, \quad M_t =  \|\hparam_t\|_{V_t^2}^2/(\|\hat{\theta}_t\|_{V_t}^2  - (\beta^{\PVT}_t)^2)\$ and $k\in[d]$ is such that $\lambda_{k}(V) \leq M_t \leq \lambda_{k+1}(V) $.
	Then for all $t \in [T]$, conditioned on event $\hat{\cA}_t \cap \tilde{\cA}_t$ , it holds for all $s \leq t$ that 	\$
	\inf _{\theta \in \cE_t^{\PVT}}\|x^\star_s(\theta)\|_{V_s^{-1}} \geq \Psi_s,  \quad   \sup _{\theta \in \cE_s^{RLS}}\|x^\star_s(\theta)\|_{V_s^{-1}} \leq \Phi_s , \quad \alpha_s\leq \halpha_s\coloneqq \Phi_s/\Psi_s.\$

\end{theorem}

\begin{comment}
	\begin{remark}
We remark that 	$ \hat{\alpha}_t$ only depends on the position (through $\hat{\theta}_t$), shape (through $V_t$), and alignment (through the relationship between $V_t$, $\hat{\theta}_t$ and $\beta_t$) of the confidence ellipsoid $\cE_t$ and can be computed efficiently. 
\end{remark}
\end{comment}

To better understand what Theorem \ref{thm:alpha} implies, we discuss some special cases. 

\paragraph{Case 1: a pure exploration regime.} \label{eg:pure} When the decision-maker doesn't care about the regret and adopts a pure exploration algorithm that plays actions in all directions sequentially, we expect $\lambda_i(V_t) = \cO(t)$ for all $i = 1,\dots, d$. Then $\hat{\alpha}_t \leq C$ for some constant $C>0$. Specifically, if $V_t = D \II_d$ for some constant $D>0$, we have $\|x\|_{V_t^{-1}} =D^{-1} \|x\| = D^{-1}$ for all $x \in \cS_{d-1}$ and hence $ \hat{\alpha}_t = 1$  .

\paragraph{Case 2:  linear structures.} 
\label{eg:principle_eigen}

When the reward takes a linear form in the action, for example, the stochastic linear bandits we study,  it's observed in practice that the estimate $\hat{\theta}_t$ tends to align in the first eigenspace of the uncertainty structure $V_t$. To see this, note that in order to maximize the reward in an online manner, bandit algorithms tend to select actions towards the confidence ellipsoid centered at  $\hat{\theta}_t$, hence forcing more exploration along the direction of $\hat{\theta}_t$, especially at the late stage of the algorithm when the ellipsoid is small. The following proposition states that, as long as the center $\hparam_t$ of the confidence ellipsoid $\cE_t$  tends to be aligned with the first eigenspace of $V_t$, the corresponding sensitivity ratio tends to 1, indicating a regret bound matching the minimax $\tilde{\cO}(d\sqrt{T})$ rate by Theorem~\ref{thm:oracle_regret}. This provides a plausible explanation for the empirical success of LinTS and Greedy.
\begin{proposition}\label{prop:align}
	 Suppose $\lim_{t\to \infty}\frac{\|\hparam_t\|_{V_{t}}^2}{\|\hparam_t\|^2} =\lambda_{\max}(V_{t}) $, it holds that $\lim_{t\to \infty} \hat{\alpha}_t = 1$.	
\end{proposition}





\subsection{TS-MR and Greedy-MR Algorithms}




In the following, we propose the algorithm Linear Thompson Sampling with Maximum (Regret) Ratio (TS-MR). The intuition is to detect the performance of LinTS by estimating $\mu_t$ and avoid bad LinTS actions by switching to OFUL occasionally. 

Specifically, at each time $t$, TS-MR calculates an estimate $\hat{\mu}_t$ and compares it with a preset threshold $\mu$. If  $\hat{\mu}_t > \mu$, LinTS might be problematic and  TS-MR takes an OFUL action to ensure a low instantaneous regret; if $\hat{\mu}_t\leq \mu$, TS-MR  takes the LinTS action. We remark that setting $\iota_t \equiv 0$ yields the corresponding Greedy-MR algorithm. The pseudocode is presented in Appendix~\ref{sec:TS-MR}.


By construction, TS-MR and Greedy-MR have $\mu_t \leq  \max\{\mu,2\}$ for all $t\in[T]$. Using Theorem~\ref{thm:oracle_regret}, we get the following frequentist regret for them, which is optimal up to a constant factor.
\begin{corollary} TS-MR and Greedy-MR achieve a frequentist regret of
$
	\cR(T) = \tilde{\cO}\left( \max\{\mu,2\} d \sqrt{T}\right)
$.
\end{corollary}









